\section{Particle Geometry} \label{New Fundamental}

The Geo model treats collision response as a current-frame geometric operation.
The response is computed from the state that exists in the current simulation
step:

\begin{itemize}[leftmargin=*]
    \item current particle positions,
    \item current particle velocities,
    \item current contact normal,
    \item current penetration depth,
    \item current overlap area,
    \item and the source particle's simultaneous contact areas.
\end{itemize}

The model should not require a remembered first-contact velocity, a stored
turnaround depth, or a persistent compression/rebound phase.  Those quantities
may be useful diagnostics, but they are not required dynamics state in the
instantaneous geometric model.

\section{Current Relative Normal Velocity}

For a source particle $s$ and target particle $t$, let the current contact
normal be:

\begin{equation}
\begin{aligned}
\hat{n}_{s,t}
=
\frac{x_t - x_s}{\lVert x_t - x_s \rVert}
\end{aligned}
\end{equation}

The current relative normal velocity is:

\begin{equation}
\begin{aligned}
v_n
=
(v_t - v_s) \cdot \hat{n}_{s,t}
\end{aligned}
\end{equation}

With this sign convention:

\begin{itemize}[leftmargin=*]
    \item $v_n < 0$ means the source and target are closing,
    \item $v_n = 0$ means there is no relative normal motion,
    \item $v_n > 0$ means the source and target are separating.
\end{itemize}

The collision response uses this current-frame value.  It is not the velocity
from the first frame of a collision episode.

\section{Overlap Geometry}

For two circular particles, the current overlap area is:

\begin{equation}
\begin{aligned}
A_{s,t} = A(r_s, r_t, d)
\end{aligned}
\end{equation}

where $r_s$ and $r_t$ are the particle radii and $d$ is the current center
distance.

The overlap area is a geometric property of the contact.  The source-side and
target-side overlap areas are the same area:

\begin{equation}
\begin{aligned}
A_{s,t} = A_{t,s}
\end{aligned}
\end{equation}

Velocities can change the response, but they do not change the instantaneous
geometric overlap area.

\section{Geometric Response}

The current-frame geometric response should be a function of overlap area,
current relative normal velocity, stiffness, and source-area weighting:

\begin{equation}
\begin{aligned}
\Delta v_{s,t}
=
G(A_{s,t}, v_n, Q, w_{s,t})
\end{aligned}
\end{equation}

where:

\begin{itemize}[leftmargin=*]
    \item $A_{s,t}$ is the current overlap area,
    \item $v_n$ is the current relative normal velocity,
    \item $Q$ is the stiffness parameter,
    \item $w_{s,t}$ is the source-owned area weight for this contact.
\end{itemize}

The exact form of $G$ is the next dynamics decision.  The important convention
is that $G$ must be computed from the current frame.  It should not depend on
stored first-contact values.

\section{Source-Owned Area Weighting}

A source particle may contact multiple targets in the same frame.  The source
particle's total current overlap area is:

\begin{equation}
\begin{aligned}
A_s^{total}
=
\sum_{t \in C_s} A_{s,t}
\end{aligned}
\end{equation}

where $C_s$ is the current contact set for source particle $s$.

The area weight for one contact is:

\begin{equation}
\begin{aligned}
w_{s,t}
=
\frac{A_{s,t}}{A_s^{total}}
\end{aligned}
\end{equation}

when $A_s^{total} > 0$.

These weights keep the source particle from responding to each contact as if it
were the only contact.  In a multi-contact frame, the source response must be
distributed across the current overlap areas.

\section{Legacy Turnaround Formulas}

Previous notes used an incoming collision velocity $v_{in}$ and a stored
turnaround depth:

\begin{equation}
\begin{aligned}
\alpha_0
=
\left(\frac{3\mu v_{in}^2}{2Q}\right)^{1/3}
\end{aligned}
\end{equation}

This formula remains useful as a diagnostic or comparison curve.  In the new
Geo model, however, $\alpha_0$ should not be stored as persistent dynamics
state.  If a zero-velocity or turnaround estimate is needed, it should be
derived from the current frame and treated as temporary diagnostic information.

\section{Starting In Contact}

When particles start in contact, the response should still be computed from the
current frame:

\begin{itemize}[leftmargin=*]
    \item current overlap area,
    \item current relative normal velocity,
    \item stiffness,
    \item and current source-area weights.
\end{itemize}

The simulation should not need to infer an earlier collision history.  If the
starting overlap is too deep for the selected geometric response rule, that is a
guardrail or configuration issue, not a reason to create hidden collision
history.
